\begin{frame}{역사: n-gram에서 신경망까지}
  \begin{itemize}
    \item ``다음 단어를 맞히기''는 현대 LLM이 발명한 것이
      아님. \kb{1940년대 섀넌}(Claude Shannon)의 정보이론
      논문에서 영어의 ``다음 단어 예측 엔트로피''를 추정할
      때 쓰던 것이 바로 n-gram — 사람에게 문장을 읽고 다음
      단어를 맞히게 해 기계의 상한을 재는 실험까지.
    \item \kb{1990$\sim$2000년대}: n-gram($n=3\sim5$)이
      음성인식·기계번역의 표준. 그러나 \textbf{두 벽}:
      (1) \textbf{문맥 길이} — 앞의 $n-1$개만 봄,
      (2) \textbf{OOV와 조합 폭발} — $V^n$($V=50{,}000, n=3$이면
      $1.25 \times 10^{14}$개) 중 대부분 0 확률 $\to$
      \kb{스무딩}(등장하지 않은 조합에도 작은 확률 분배)
      필수.
    \item \kb{2013년 word2vec}(Mikolov et al.): ``단어의 의미 =
      주변 단어 분포''라는 분포 가설을 \textbf{임베딩}으로
      효율 학습(Week 14).
    \item \kb{2017년 Transformer}(Week 12): 자기-어텐션으로
      ``문장 전체의 모든 토큰을 동시에 조건부''로 보고,
      \textbf{스케일만} 키워도 성능이 계속 올라감 $\to$ 80년
      된 과제가 \kb{LLM}이라는 새로운 현상(Brown et al.,
      2020)을 만들어냄.
  \end{itemize}
\end{frame}
